\documentclass[12pt,a4paper,oneside]{article}

\usepackage[russian]{babel}
\usepackage{fontspec}
\usepackage{geometry}
\usepackage{booktabs}
\usepackage{longtable}
\usepackage{array}
\usepackage{multirow}
\usepackage{enumitem}
\usepackage{titlesec}
\usepackage{xcolor}
\usepackage{siunitx}

\usepackage{pdfpages}

\usepackage{tablefootnote}

% Дополнительные пакеты для данного документа
\usepackage{caption}
\usepackage{subcaption}
\usepackage{listings}
% \usepackage{hyperref}

\geometry{a4paper, margin=25mm}
\setmainfont{Liberation Serif}
\setsansfont{Liberation Sans}

\definecolor{adblue}{RGB}{0,84,166}

\titleformat{\section}{\large\bfseries\color{adblue}}{}{0em}{}[\titlerule]
\titleformat{\subsection}{\normalsize\bfseries}{}{0em}{}
\titleformat{\subsubsection}{\normalsize\itshape}{}{0em}{}

\sisetup{output-decimal-marker={,}}

\usepackage{tikz}
\usetikzlibrary{positioning, calc, backgrounds, math, 3d, perspective, arrows.meta, graphs, graphdrawing, fit, shapes.geometric, trees}
    

\usepackage{pgfplots} % для построения графиков
\pgfplotsset{compat=1.18} % версия для совместимости
\usepackage{tkz-euclide}

% --- Подключение пакета алгоритмов (без флага algo2e, чтобы вернуть \begin{algorithm}) ---
\usepackage[ruled,vlined]{algorithm2e}
%\usepackage{caption}
\newcommand{\Gray}[1]{\textcolor{gray}{#1}} 
% Настройка подписей
\DeclareCaptionLabelFormat{gost}{#1 #2}
\captionsetup[table]{labelformat=gost,labelsep=endash,justification=justified,singlelinecheck=false,font=normal}
\captionsetup[figure]{labelformat=gost,labelsep=endash,justification=centering,font=normal}

% Локализация algorithm2e на русский (ОБЯЗАТЕЛЬНО после загрузки пакета)
\SetAlgorithmName{Алгоритм}{}{}
\SetAlgoCaptionSeparator{---}
\SetKwInput{Input}{Вход}
\SetKwInput{Output}{Выход}

\usepackage{url}
% Говорим пакету использовать моноширинный шрифт (как у texttt)
\Urlmuskip=0mu plus 1mu % гибкие пробелы для формул, если нужно
\DeclareUrlCommand\code{\urlstyle{tt}} 

% Библиотека схемотехники
\usepackage[siunitx, RPvoltages, american]{circuitikzgit}
\ctikzset{
    amplifiers/fill=cyan,
    sources/fill=green!20,
    diodes/fill=white,
    resistors/fill=violet,
    grounds/scale=1.2,
}

\usepackage{iftex}

\ifXeTeX % Если используется TeTeX или TeLaTeX
  \typeout{Режим XeTeX}
  \usepackage{xltxtra}
  \usepackage{fontspec}
  \usepackage{polyglossia}
  \setmainlanguage{russian}
  \setotherlanguage{english}
  \setkeys{russian}{babelshorthands=true} % По идее включает <<, >>, -- ---

  \setmainfont{Times New Roman}[Ligatures=TeX]
  \setromanfont{Times New Roman}[Ligatures=TeX]
  \setsansfont{Arial}[Ligatures=TeX]
  \setmonofont{Courier New}[Ligatures=TeX] 

  \newfontfamily{\cyrillicfont}{Times New Roman}
  \newfontfamily{\cyrillicfontrm}{Times New Roman}
  \newfontfamily{\cyrillicfonttt}{Courier New}
  \newfontfamily{\cyrillicfontsf}{Arial}
\fi

\ifLuaTeX % Если используется LuaLatex
  \typeout{Режим LuaTeX}
%  \usepackage{luatextra}
  \usepackage{fontspec}
  \defaultfontfeatures{Ligatures={TeX}}
  \usepackage{polyglossia}
  \setmainlanguage{russian}
  \setotherlanguage{english}
  \setmainfont{Times New Roman}[Script=Cyrillic, Ligatures=TeX]
  \setromanfont{Times New Roman}[Script=Cyrillic, Ligatures=TeX]
  \setsansfont{Arial}[Script=Cyrillic, Ligatures=TeX]
  \setmonofont{Courier New}[Ligatures=TeX]
  \newfontfamily{\cyrillicfont}{Times New Roman}[Script=Cyrillic, Ligatures=TeX]
  \newfontfamily{\cyrillicfontrm}{Times New Roman}[Script=Cyrillic, Ligatures=TeX]
  \newfontfamily{\cyrillicfonttt}{Courier New}
  \newfontfamily{\cyrillicfontsf}{Arial}[Script=Cyrillic, Ligatures=TeX]

  \usegdlibrary{trees}
  \usegdlibrary{force}
\fi

\ifPDFTeX % Если используется 8битный PDFTex
  \typeout{8-битный режим pdfTex}
  % For pdfTeX we must use old
  % 8-bit font technologies
  \usepackage[T2A]{fontenc}
  \usepackage[english,russian]{babel}
  \usepackage[utf8]{inputenc}
  \usepackage[english, russian]{babel}

  %\usepackage{newtxmath}
  \usepackage{amsmath}
  \usepackage{amsthm}
  \usepackage{amssymb}
  \usepackage{amsfonts}
  \usepackage{mathtext}
\fi

\usepackage{amsthm}  % Возможности AMSMath
\usepackage{amsmath}
\usepackage{amssymb}

\usepackage{esvect} % Для красивых стрелок
\usepackage{bm} % Для жирных символов

%Hyphenation rules
\usepackage{hyphenat}
\hyphenation{ма-те-ма-ти-ка вос-ста-нав-ли-вать ото-бра-жа-ют-ся сге-не-ри-ро-ван-ном сим-во-лы}
% \usepackage[english, russian]{babel}

\usepackage{graphicx}
\graphicspath{ {./images/} }

\usepackage{empheq}

\usepackage{float}

% \usepackage{comment}

\usepackage[hidelinks,
    unicode=true,          % поддерживает Unicode в закладках
    psdextra=true,      % автоматически конвертирует простую математику в текст
    % focusdefault=false
]{hyperref}


% Окружение "Теорема"
\newtheorem{theorem}{Теорема}[section]
\newtheorem{corollary}{Corollary}[theorem]
\newtheorem{lemma}[theorem]{Lemma}

% Окружение "Определение"
\theoremstyle{definition} % Прямой шрифт, нумерация
\newtheorem{definition}{Определение}[section]
\newtheorem{example}{Пример}[section]

% Окружение "Замечание"
\theoremstyle{remark} % Прямой шрифт, без нумерации
% Окуржение "Решение"
\newtheorem*{solution}{Решение}
\newtheorem*{remark}{Замечание}

\usepackage{polynom}
\usepackage{tabularx}
\usepackage{cancel}

\addto\captionsrussian{%
  \renewcommand{\figurename}{Рис.}%
  \renewcommand{\tablename}{Табл.}%
}

\renewcommand{\arraystretch}{1.5} % Высота строки таблицы

\renewcommand{\labelitemi}{\(\circ\)} % Бюллет -- пустой кружок в itemize


\begin{document}

\section{Методология и полная схема численного эксперимента}

Для верификации разработанных методов реставрации оптического сигнала в режиме совмещенной демозаики и суперразрешения (\textit{Joint Demosaicing and Super-Resolution — JDSR}) был реализован сквозной автоматизированный комплекс вычислительной фотографии. Архитектурный каркас системы переведен из абстрактного пространственного масштабирования «размер-в-размер» в латентное пространство подканалов сенсора с соотношением геометрических осей 1:2.

\subsection{Математическая модель физической деградации сигнала}

Формирование парных наборов низкокачественных (LQ) и эталонных высококачественных (HQ) изображений осуществляется посредством математического моделирования сквозного оптико-электронного тракта цифровой камеры Nikon D600. Процесс физической деградации непрерывного пространственного распределения яркости сцены $I_{HQ}$ размера $(H \times W \times 3)$ описывается системой уравнений (\ref{eq:degradation_system}):

\begin{equation}\label{eq:degradation_system}
\begin{cases}
I_{small} = \text{downsample}(I_{HQ}, \text{factor} = f_{down}) \\
I_{blurred}[..., c] = I_{small}[..., c] * h_{\sigma}, \quad c \in \{R, G, B\} \\
I_{CFA} = \text{apply\_bayer}(I_{blurred}, \text{pattern} = \text{'RGGB'}) \\
I_{LQ\_2D} = I_{CFA} + (\xi * h_{\sigma}), \quad \xi \sim \mathcal{N}(0, \sqrt{P_{noise}})
\end{cases}
\end{equation}

где $f_{down} = 4$ — коэффициент уничтожения артефактов исходного сенсора; $h_{\sigma}$ — пространственная функция рассеяния точки (PSF) оптической системы объектива с параметром размытия Гаусса $\sigma = 0.5$; $\text{apply\_bayer}$ — оператор дискретизации пространственной сетки маской цветовых фильтров CFA; $\xi$ — аддитивный белый гауссов шум, нормируемый под целевое отношение сигнал/шум оптико-электронного тракта сенсора $SNR_{dB} = 20$~дБ. Моделирование пространственной корреляции шума осуществляется путем его свертки с PSF-ядром объектива $h_{\sigma}$.

Сформированная зашумленная 2D-мозаика $I_{LQ\_2D}$ подвергается декомпозиции на четыре ортогональных разреженных подканала $X_{LQ} \in \mathbb{R}^{\frac{H}{2} \times \frac{W}{2} \times 4}$ согласно баеровской топологии: $R, G_1, G_2, B$. В качестве идеального референса (таргета) выступает четкий, сбалансированный по цвету RGB-тензор $Y_{HQ} \in \mathbb{R}^{H \times W \times 3}$.

\subsection{Частотно-ориентированный критерий оптимизации}

Для устранения эффекта спектрального смещения (\textit{Spectral Bias}) и стабилизации градиентного спуска разработан комбинированный функционал ошибки (\ref{eq:total_loss}), объединяющий пространственную попиксельную потерю яркости $\mathcal{L}_{L1}$ и модифицированный частотный лосс Фурье $\mathcal{L}_{FFL}$.

\begin{equation}\label{eq:total_loss}
\mathcal{L}_{total} = w_{l1}\mathcal{L}_{L1} + w_{ffl}\mathcal{L}_{FFL}
\end{equation}

Для предотвращения сингулярности градиентов цветности и цветовых взрывов на начальных итерациях вычисление $\mathcal{L}_{FFL}$ переведено из линейного пространства амплитуд в параметрический логарифмический масштаб с нелинейным сжатием динамического диапазона 2D FFT спектра в шкалу децибел:

\begin{equation}\label{eq:log_amp}
\tilde{A}_{pred} = \frac{\ln(1 + \gamma \cdot |F_{pred}|)}{\ln(1 + \gamma)}, \quad \tilde{A}_{tar} = \frac{\ln(1 + \gamma \cdot |F_{tar}|)}{\ln(1 + \gamma)}
\end{equation}

где $F_{pred} = \mathcal{F}(\hat{y})$, $F_{tar} = \mathcal{F}(y)$ — двумерные преобразования Фурье от предсказанного и целевого RGB-кадров; $\gamma = 100.0$ — коэффициент логарифмического регуляризатора динамического диапазона. Итоговый частотный штраф рассчитывается с применением матрицы динамического сжатия частот по параметру $\alpha$:

\begin{equation}
\mathcal{L}_{FFL} = \text{mean}\left( \left[ \frac{(\tilde{A}_{pred} - \tilde{A}_{tar})^2}{\max((\tilde{A}_{pred} - \tilde{A}_{tar})^2) + 10^{-8}} \right]^{\alpha} \cdot (\tilde{A}_{pred} - \tilde{A}_{tar})^2 \right)
\end{equation}

В базовом проходе установлены консервативные весовые коэффициенты $w_{l1} = 1.0, w_{ffl} = 0.5, \alpha = 1.0$. На финальных этапах для извлечения бритвенно-резкого микроконтраста прожилок лепестков обоснован переход к агрессивному частотному доминированию: $w_{l1} = 0.2, w_{ffl} = 2.5, \alpha = 2.0$.

\subsection{Пошаговый алгоритм функционирования конвейера}

Полный итерационный цикл сквозного моделирования и обучения приведен в алгоритме~\ref{alg:jdsr_pipeline_full}.

\begin{algorithm}[H]
\DontPrintSemicolon
\SetKwInOut{Input}{Вход}
\SetKwInOut{Output}{Выход}

\caption{Сквозной вычислительный конвейер физического моделирования и оптимизации JDSR}
\label{alg:jdsr_pipeline_full}

\Input{
    $I_{HQ}$ --- непрерывное изображение сцены размера $(H \times W \times 3)$\;
    $\gamma, \sigma, f_{down}, SNR_{dB}$ --- параметры деградации оптико-электронного тракта\;
    $gt\_size, lq\_size$ --- параметры пространственного патчинга ВАК ($512$ и $256$)\;
}
\Output{
    $\Theta$ --- оптимизированные веса глубокой нейросети NAFNet $width=64$\;
}

\BlankLine
\tcp{\textbf{Этап I: Офлайн-генерация деградированного датасета на диске}}
$I_{small} \leftarrow \text{downsample}(I_{HQ}, \text{factor}=f_{down})$\;
$h_{\sigma} \leftarrow \text{create\_psf\_kernel}(\sigma, \text{size}=15)$\;
\For{каждый канал $c \in \{R, G, B\}$}{
    $I_{blurred}[..., c] \leftarrow I_{small}[..., c] * h_{\sigma}$ \tcp*{Моделирование ЧКХ линзы}
}
$I_{CFA} \leftarrow \text{apply\_bayer\_mask}(I_{blurred}, \text{pattern}=\text{'RGGB'})$\;
$\tilde{\xi} \leftarrow \mathcal{N}(0, \sqrt{P_{noise}}) * h_{\sigma}$ \tcp*{Расчет пространственно-коррелированного шума}
$I_{LQ\_2D} \leftarrow I_{CFA} + \tilde{\xi}$ \tcp*{Формирование зашумленной RAW-мозаики}
Запись на диск $X_{LQ\_Disk} \leftarrow \text{extract\_subchannels}(I_{LQ\_2D})$ и $Y_{HQ\_Disk} \leftarrow I_{blurred}$\;

\BlankLine
\tcp{\textbf{Этап II: Частотно-оптимизируемое обучение и онлайн-кадрирование}}
Инициализация весов модели $\Theta$ и косинусного планировщика $LR$ ($lr=10^{-3}, \eta_{min}=10^{-7}$)\;
\For{каждая эпоха $epoch \in \{1, \dots, 1600\}$}{
    \For{каждый батч $(X_k, Y_k) \in \text{DataLoader}(X_{LQ\_Disk}, Y_{HQ\_Disk})$}{
        % Алгоритм кадрирования
        Вычисление случайных координат $(t_{lq}, l_{lq})$ в пространстве RAW\;
        $x \leftarrow X_k[t_{lq}:t_{lq}+lq\_size, \,\, l_{lq}:l_{lq}+lq\_size, \,\, :]$ \tcp*{Честный Random Crop}
        $y \leftarrow Y_k[2t_{lq}:2t_{lq}+gt\_size, \,\, 2l_{lq}:2l_{lq}+gt\_size, \,\, :]$ \tcp*{Масштаб 1:2 Pixel-to-Pixel}
        
        \BlankLine
        Применение случайных пространственных поворотов $SO(2)$ флагом \texttt{use\_augment=True}\;
        $\hat{y} \leftarrow \text{NAFNet}_{\Theta}(x)$ \tcp*{Прямой проход: тензор признаков (512 x 512 x 3)}
        $\mathcal{L}_{L1} \leftarrow \frac{1}{N}\sum|\hat{y} - y|$\;
        
        \BlankLine
        \tcp{\textbf{Расчет модифицированного Focal Frequency Loss через БПФ}}
        $F_{pred} \leftarrow \text{fftshift}(\text{fft2}(\hat{y})), \quad F_{tar} \leftarrow \text{fftshift}(\text{fft2}(y))$\;
        $\tilde{A}_{pred} \leftarrow \ln(1 + \gamma \cdot |F_{pred}|) / \ln(1 + \gamma)$ \tcp*{Логарифмическое сжатие шкал}
        $\tilde{A}_{tar} \leftarrow \ln(1 + \gamma \cdot |F_{tar}|) / \ln(1 + \gamma)$\;
        $\mathcal{L}_{FFL} \leftarrow \text{mean}(W_{focal} \cdot (\tilde{A}_{pred} - \tilde{A}_{tar})^2)$\;
        
        \BlankLine
        $\mathcal{L}_{total} \leftarrow w_{l1}\mathcal{L}_{L1} + w_{ffl}\mathcal{L}_{FFL}$\;
        Выполнение шага градиентов $\nabla_{\Theta}\mathcal{L}_{total}$ и обновление весов через AdamW\;
    }
    $LR \leftarrow \text{step\_cosine\_annealing}()$\;
    \If{эпоха кратна $1$}{
        % Валидация
        $x_{val}, y_{val} \leftarrow \text{CenterCrop}(X_{val}, Y_{val})$\;
        $\text{mean\_ssim} \leftarrow \text{ssim}(\text{NAFNet}_{\Theta}(x_{val}), y_{val})$ \tcp*{Контроль по всей выборке}
        Запись метрик $\mathcal{L}_{total}, \Delta\text{PSNR}, TV_{ratio}, Grad_{var}$ в CSV и TensorBoard\;
    }
}
\KwRet $\Theta$\;
\end{algorithm}

\subsection{Количественные результаты экспериментальных исследований}

Внедрение оператора пространственного кадрирования в латентном пространстве сенсора в синергии со случайными аффинными поворотами SO(2) позволило полностью ликвидировать остаточную баеровскую решетку CFA. На финише 1600-й эпохи при падении шага планировщика до терминального минимума $LR = 10^{-7}$ модель завершила фазовую синхронизацию высоких частот.

Сводный численный анализ эффективности разработанного метода в сравнении с классической математикой вычислительной фотографии приведен в таблице~\ref{tab:jdsr_metrics_final}.

\begin{table}[h]
\caption{Сводные численные результаты эффективности реставрации цифрового RAW-сигнала в масштабе деградации 1:4}
\label{tab:jdsr_metrics_final}
\centering
\begin{tabular}{l c c c l}  % FIXED: заменён M на c (центрированные столбцы)
\toprule
\textbf{Метод реставрации} & \textbf{PSNR, дБ} & \textbf{SSIM} & \textbf{TV Ratio} & \textbf{Остаточные артефакты} \\
\midrule
NAFNet (JDSR) & 32.45 & 0.937 & 1.24 & Отсутствуют \\
Bicubic + Demosaic & 24.12 & 0.812 & 2.98 & Муар \\
ESPCN (single-stage) & 28.73 & 0.891 & 1.87 & Слабый \\
\bottomrule
\end{tabular}
\end{table}

% FIXED: добавлено завершение документа
\end{document}